%---------------------------------------------------------------------------
\chapter{線形システム論（離散系）}
%---------------------------------------------------------------------------

離散システムは連続システムを離散化してから制御系を設計する場合や，もともと離散システムで表されるシステム（計算機内での計算アルゴリズムや，経済システム（四半期ごとの指標をモデル化する場合など））を制御する場合に用いられる．

離散システムの解析設計手法は，連続時間システムの場合とほとんど同様に議論できるので，ここでは連続時間システムに対する議論との対比を行い，連続時間システム
の場合と異なる部分について述べる．
%---------------------------------------------------------------------------
\section{状態方程式}
%---------------------------------------------------------------------------
本章では以下の状態方程式で表されている離散システムを扱うものとする．
\begin{align}
 x_{k+1} &= A x_k + B u_k \label{d_state_eq} \\
 y_k     &= C x_k + D u_k \label{d_output_eq}
\end{align}
ここで$x_k \in \mathbb{R}^n$, $u_k \in \mathbb{R}^m $, $y_k \in
\mathbb{R}^p $はそれぞれ$k$時点での状態，入力，出力である．
%---------------------------------------------------------------------------
\section{座標変換}
%---------------------------------------------------------------------------
システムの座標変換も連続時間システムと同様に定義できる．
状態$x_k$を次のような座標変換で$\bar{x}_k$に変換することを考える．
\begin{equation}
 x_k = T \bar{x}_k
\end{equation}
ここで$T$は正則な行列であるとする．1時刻後の状態$\bar{x}_{k+1}$は次のよ
うになる．
\begin{align}
 \bar{x}_{k+1} &= T^{-1} x_{k+1} \nonumber \\
               &= T^{-1} ( A x_k + B u_k) \nonumber \\
               &= T^{-1} A T \bar{x}_k + T^{-1} B u_k \nonumber \\
               &\defeq \bar{A} \bar{x}_k + \bar{B} u_k
\end{align}
これが状態を$\bar{x}_{k+1}$としたときの状態方程式となる．また出力方程式は
\begin{align}
 y_k &= C T \bar{x}_k + D u_k \nonumber \\
     &\defeq \bar{C} \bar{x}_k + \bar{D} u_k
\end{align}
となり，座標変換による$A$, $B$, $C$, $D$行列の変化は連続時間システムの場
合と一致する．
%---------------------------------------------------------------------------
\section{伝達関数と状態方程式}
%---------------------------------------------------------------------------
連続時間システムを伝達関数表現するためにはシステムのラプラス変換を用いた．
離散時間システムに対しても同様に$z$変換というものがあり，伝達関数表現す
ることがある．しかし，ここでは直接コントローラの設計に用いないので説明は
省略する．
%---------------------------------------------------------------------------
\section{可制御性，可観測性}
%---------------------------------------------------------------------------
連続時間システムと離散時間システムの状態の座標変換に対する状態方程式の各
行列の変換は厳密に一致している．よって厳密さを求めなければ連続時間システ
ムで行った議論がそのまま離散時間システムにも成り立つ．
%---------------------------------------------------------------------------
\section{システムの応答}
%---------------------------------------------------------------------------
状態方程式(\ref{d_state_eq})より明らかに
\begin{align*}
 x_1 &= A x_0 + B u_0 \\
 x_2 &= A x_1 + B u_1 \\
     &= A (A x_0 + B u_0) + B u_1 \\
     &= A^2 x_0 + A B u_0 + B u_1 \\
 x_3 &= A x_2 + B u_2 \\
     &= A (A^2 x_0 + A B u_0 + B u_1) + B u_2 \\
     &= A^3 x_0 + A^2 B u_0 + A B u_1 + B u_2
\end{align*}
となる．一般に
\begin{equation}
 x_k = A^k x_0 + \sum_{i=0}^{k-1} A^{k-i-1} B u_i
\end{equation}
である．入出力応答は
\begin{equation}
 y_k = C A^k x_0 + \sum_{i=0}^{k-1} C A^{k-i-1} B  u_i + D u_k
\end{equation}
となる．
%---------------------------------------------------------------------------
\section{システムの安定性}
%---------------------------------------------------------------------------
いまシステムの入力が$u_i = 0$の場合を考える．このときのシステムの応答は
\begin{equation}
 x_k = A^k x_0
\end{equation}
となる．このとき，
\begin{equation}
 \lim_{k \rightarrow \infty} A^k = 0
\end{equation}
であればシステムの状態も
\begin{equation}
 \lim_{k \rightarrow \infty} x_k = 0
\end{equation}
となり，システムが安定であるといえる．さて正則な複素行列$T$を
\begin{align}
 T^{-1} A T
  &=
  \left(\begin{array}{cccc}
         \lambda_1 & 0         & \cdots & 0 \\
	 0         & \lambda_2 &        & 0 \\
	 \vdots    &           & \ddots & \\
	 0         & 0         &        & \lambda_n 
	\end{array}\right) \nonumber \\
 &\defeq \Lambda
\end{align}
を満たすように選べると仮定する(ジョルダン形式になる場合も後の議論はほと
んど同じになる)．ここで$\lambda_i$は行列$A$の固有値である．この
$\Lambda$を用いれば
\begin{align}
 A^k &= (T \Lambda T^{-1})^k = T \Lambda ^k T^{-1} \nonumber \\
 &= T
  \left(\begin{array}{cccc}
         \lambda_1^k & 0           & \cdots & 0 \\
	 0           & \lambda_2^k &        & 0 \\
	 \vdots      &             & \ddots & \\
	 0           & 0           &        & \lambda_n^k
	\end{array}\right)
  T^{-1}
\end{align}
となる．よって
\begin{equation}
 \lim_{k \rightarrow \infty} A^k = 0
\end{equation}
\[\Updownarrow\]
\begin{equation}
 \lim_{k \rightarrow \infty} \lambda_i^k = 0, \quad
  i = 1, 2, \cdots, n
\end{equation}
\[\Updownarrow\]
\begin{equation}
 \vert \lambda_i \vert < 1, \quad i = 1, 2, \cdots, n
\end{equation}
であることが分かる．つまり，行列$A$の固有値の絶対値が1未満であればシステ
ムは安定であることが分かる．
%---------------------------------------------------------------------------
\section{状態フィードバック}
%---------------------------------------------------------------------------
これは連続時間システムとまったく同様に考えることができる．いま状態$x_k$が
利用できるとする．このとき入力を
\begin{equation}
 u_k = F x_k + v_k
\end{equation}
とすれば状態方程式は
\begin{equation}
 x_{k+1} = (A + B F)x_k + B v_k
\end{equation}
となる．ここで$v_k \in \mathbb{R}^m $は新しい入力である．このように状態
フィードバックされた閉ループ系の安定性は$(A + BF)$の固有値によって決まる．
これは連続時間システムの場合とまったく同じである．状態フィードバックによ
りシステムを安定化したり，システムの極を変えたりすることができる．
%---------------------------------------------------------------------------
\subsection{極配置問題}
%---------------------------------------------------------------------------
離散時間システムの場合も連続時間システムの場合と同じく，$(A+BF)$の固有値を
任意の与えられた値にする問題を極配置問題という．連続時間システムと離散時
間システムの座標変換がまったく同じ形で表されることに注意すれば，離散時間シ
ステムに対する極配置問題は連続時間システムと同一の手法で解けることがわか
る．ただし，ここではシステムの安定性を保証するために指定する極は安定極
(絶対値が1未満)である必要がある．
%---------------------------------------------------------------------------
\subsection{最適制御}
%---------------------------------------------------------------------------
安定化フィードバックとして，次の2次形式評価関数を最小にするように選ぶ方法(最
適フィードバック)がある．
\begin{equation}
 J = \sum_{i=0}^{\infty} x_i^T Q x_i + u_i^T R u_i 
\end{equation}
ここで$Q$は半正定($Q = H^T H$で表わされる)，$R$は正定($R = L^T L$,
($L$は正則行列)で表わされる)とする．この$J$を最小にするフィードバッ
クは$P$をRiccati方程式
\begin{align}
 P   &= A^T P A + Q - A^T P B (R + B^T P B)^{-1} B^T P A
\end{align}
の正定解として
\begin{align}
 u_k &= - (R + B^T P B)^{-1} B^T P A x_k \defeq Fx_k
\end{align}
で表わされることが知られている．さらに$(H, A)$が可観測であれば閉ループ系
が安定になる．
一般に$Q$と$R$の選び方は試行錯誤となるが，選び方の指針は連続時間システム
と同じである．

Riccati方程式の解法としては数値的に安定な以下の方法がよく用いられる．い
ま$A$が正則であるとする(一般に連続時間システムを離散化したとき，離散時間
システムの$A$は正則になる)．行列$\Phi$を次のように定義する．
\begin{equation}
 \Phi =
  \left(\begin{array}{cc}
         A + B R^{-1} B^T A^{-T} Q & - B R^{-1} B^T A^{-T}\\
	 - A^{-T} Q                &  A^{-T}
	\end{array}\right)
\end{equation}
この$\Phi$の固有値のうち絶対値が1未満なものが$n$個あることが分かっている．
この$n$個の固有値に対応する固有ベクトルを$w_1, w_2, \cdots, w_n$として，
次のように分解されているとする．
\begin{equation}
 w_i = \left(\begin{array}{c}
	      \xi_i \\
	      \eta_i
	     \end{array}\right)
\end{equation}
このとき
\begin{equation}
 P = (\eta_1, \eta_2,\cdots,\eta_n)(\xi_1,\xi_2,\cdots,\xi_n)^{-1}
\end{equation}
と選べば$P$は正定かつRiccati方程式を満たすことが知られている．
%---------------------------------------------------------------------------
\subsection{サーボ系}
%---------------------------------------------------------------------------
ここでは特にステップ目標値($r$が一定値)の場合について扱う．連続時間シ
ステムの場合と同様に積分器を初めに導入する方法を示す．
まず，システム(\ref{d_state_eq})(\ref{d_output_eq})において$D=0$とし
た以下の系を考える．
\begin{align}
   x_{k+1} &= A x_k + B u_k \label{d_st_proper_state_eq} \\
       y_k &= C x_k
\label{d_st_proper_output_eq}
\end{align}
ここで$y_k \in \mathbb{R}^m$であり，入力と出力の次数が等しいとする．

積分器に対応するものは
\begin{equation}
 \eta_{k} = \sum_{i=1}^{k} (r_{i-1}-y_{i-1}) 
\end{equation}
である．1サンプル時間後の$\eta_{k+1}$は
\begin{equation}
 \eta_{k+1} = \sum_{i=1}^{k+1} (r_{i-1}-y_{i-1}) 
   = \sum_{i=1}^{k} (r_{i-1}-y_{i-1}) + (r_k-y_k) = \eta_k + (r_k - y_k) 
\end{equation}
と差分方程式で表される．この積分器とシステム(\ref{d_st_proper_state_eq})
(\ref{d_st_proper_output_eq})を合わしたシステム(拡大系)は
\begin{align}
 \left(\begin{array}{c}
	x_{k+1} \\
	\eta_{k+1}
       \end{array}\right)
 &=
 \left(\begin{array}{cc}
        A  & 0 \\
	-C & I
       \end{array}\right)
 \left(\begin{array}{c}
        x_k \\
	\eta_k
       \end{array}\right)
 +
 \left(\begin{array}{c}
	B \\
	0
       \end{array}\right)
 u_k +
 \left(\begin{array}{c}
        0 \\
	I
       \end{array}\right) 
 r_k \\
 y_k &=
  \left(\begin{array}{cc}
         C & 0 
	\end{array}\right)
  \left(\begin{array}{c}
         x_k \\
	 \eta_k
	\end{array}\right)
\end{align}
となる．このシステムは$(A, B)$可制御かつ
\begin{equation}
 \rank \left(\begin{array}{cc}
	      A - I  & B \\
	      C      & 0
	     \end{array}\right)
 = n + m
\end{equation}
のとき安定化できることが知られている．このとき，この拡大システムを安定化さ
せる状態フィードバックは極配置法や最適制御法で求めることができる．安定化
フィードバックが
\begin{equation}
 u_k = \left(\begin{array}{cc}
	      F & K 
	     \end{array}\right)
 \left(\begin{array}{c}
        x_k \\
	\eta_k
       \end{array}\right)
\end{equation}
で与えられているとする．このとき閉ループ系は
\begin{align}
 \left(\begin{array}{c}
	x_{k+1} \\
	\eta_{k+1}
       \end{array}\right)
 &=
 \left(\begin{array}{cc}
        A + B F & B K  \\
	-C      & I
       \end{array}\right)
 \left(\begin{array}{c}
        x_k \\
	\eta_k
       \end{array}\right)
 +
 \left(\begin{array}{c}
        0 \\
	I
       \end{array}\right) 
 r_k \\
 y_k &=
  \left(\begin{array}{cc}
         C & 0 
	\end{array}\right)
  \left(\begin{array}{c}
         x_k \\
	 \eta_k
	\end{array}\right)
\end{align}
となる．

$r_k$が一定値$r$であるとすると，システムは安定であるから$i
\rightarrow \infty $で状態$x_i$と積分器の状態$\eta_i$はある値$x_{\infty}
$と$\eta_{\infty}$に収束する．よって出力$y$も$y_{\infty} = C x_{\infty}
$に収束する．$i \rightarrow \infty$において$\eta_i$が収束することから
\begin{equation}
 \lim_{i \rightarrow \infty} \{\eta_{i+1}-\eta_{i}\} = 0
\end{equation}
でなければならない．これは
\begin{align}
 \lim_{i \rightarrow \infty} \{ \eta_{i+1} - \eta_i \}
  &= \lim_{i \rightarrow \infty}
  \{ \eta_i + r - y_i - \eta_i \}\nonumber \\
 &= r - y_{\infty}\nonumber \\
 &= 0
\end{align}
であることを示している．よって定常偏差は零，つまり
\begin{equation}
 \lim_{i \rightarrow \infty} y_i = r
\end{equation}
となる．
%---------------------------------------------------------------------------
\section{出力フィードバック}
%---------------------------------------------------------------------------
これまではシステムを状態フィードバックにより安定化することを考えてきた．
しかし，一般にシステムの状態を直接知ることはできない．そこで状態観測器を
用いたフィードバック制御が必要となる．

本節ではシステム(\ref{d_st_proper_state_eq})
(\ref{d_st_proper_output_eq})を考える．状態観測器の考え方は連続時間シス
テムの場合と同一である．
%---------------------------------------------------------------------------
\subsection{同次元観測器（同一次元オブザーバ）}
%---------------------------------------------------------------------------
次のシステムを考える．
\begin{equation}
 \hat{x}_{k+1} = A \hat{x}_k + B u_k - K (y_k - C \hat{x}_k)
  \label{d_obs_org}
\end{equation}
いまシステム(\ref{d_st_proper_state_eq})の状態$x_k$とシステム
(\ref{d_obs_org})の状態$\hat{x}_k$の差を
\begin{equation}
 e_k \defeq x_k - \hat{x}_k
\end{equation}
とすれば
\begin{align}
 e_{k+1} &= x_{k+1} - \hat{x}_{k+1} \nonumber \\
         &= A x_k + B u_k 
	  - A \hat{x}_k - B u_k + K ( C x_k - C \hat{x}_k) \nonumber \\
         &= (A + KC) e_k
\end{align}
となり，$(A + KC)$が安定になるように$K$を選べば入力や初期値$\hat{x}_0$によ
らず$\hat{x}_i \rightarrow x_i$となることが分かる．つまり，連続系と同じ議論が可能となる．
%---------------------------------------------------------------------------
\subsection{状態観測器によるフィードバック系の安定性}
%---------------------------------------------------------------------------
離散時間システムにおいて状態観測器を用いた場合のフィードバック系の安定性
は，連続時間システムの場合と同様の議論により，
全体の系の極が状態フィードバックを施したときのシステ
ムの極($(A+BF)$の固有値)と状態観測器の極($(A + KC)$の固有値)で表わさ
れることが分かる．
